\begin{lemma}
The approximately normalised state given by performing QPE with $l$-ancilla qubits initialised into the $m$-th B-Spline window state is:
\begin{equation}
    \sum_i c_i \sum_{x=-2^{l-1}+1}^{2^{l-1}} \sqrt{\frac{\pi}{6m}}\left(\frac{\sin \frac{(x-E_i)\pi}{m}}{ \frac{(x-E_i)\pi}{m}}\right)^m \ket{x}\ket{\psi_i}
\end{equation}
where $E_i$ is the energy of the $i$-th eigenstate $\ket{\psi_i}$.
\label{lemma:bspline_def}
\end{lemma}
\begin{proof}
    The fourier transform of the $m$-th B-spline is given by \cite{bspline}: \begin{equation}
        \left(\frac{\sin \frac{x\pi}{m}}{ \frac{x\pi}{m}}\right)^m
    \end{equation} 
    Performing QPE with the ancillas initially in the $m$-th B-spline state is equivalent to multiplying the B-spline state by $e^{ixE_i}$ and then applying QFT. Multiplying by a phase before QFT is equivalent to shifting by that phase after QFT, hence we get amplitudes proportional to:
    \begin{equation}
     \left(\frac{\sin \frac{(x-E_i)\pi}{m}}{ \frac{(x-E_i)\pi}{m}}\right)^m 
\end{equation}
    So all we need to do is find the approximate normalisation. To find this we approximate the centre of the distribution by a Gaussian by Taylor expanding its logarithm:
    \begin{align}
                \log \left(\frac{\sin \frac{(x-E_i)\pi}{m}}{ \frac{(x-E_i)\pi}{m}}\right)^m &= m \log \frac{\sin \frac{(x-E_i)\pi}{m}}{ \frac{(x-E_i)\pi}{m}}\\
                &= m \log \left( 1 - \frac{1}{3!}(\frac{x\pi}{m})^2 + O(x^4) \right)\\ 
                &= m(- \frac{1}{6}(\frac{x\pi}{m})^2 + O(x^4))\\
                &= - \frac{x^2 \pi^2}{6m} + O(x^4)
    \end{align}
    Therefore, we can approximate the centre of the distribution by $e^{ - \frac{x^2 \pi^2}{6m}}$ and so we can utilise the normalisation for the gaussian $\sqrt{\frac{\pi}{6m}}$ to approximate the normalisation for our distribution.

\end{proof}
\begin{lemma}
    Let $\delta$ be the probability of measuring an energy $x$ outside of the confidence interval of $E_i\pm m$ when performing $m$-th B-spline boosted QPE. Then,
    \begin{equation}
        \delta = O\left(\exp (-m)\right)
    \end{equation}
    \label{lemma:bspline_prob}
\end{lemma}
\begin{proof}
From Lemma~\ref{lemma:bspline_def} we get:
\begin{equation}
     \delta = \sum_{|x - E_i| > m} \left\| \sqrt{\frac{\pi}{6m}}\left(\frac{\sin \frac{(x-E_i)\pi}{m}}{ \frac{(x-E_i)\pi}{m}}\right)^m \right\|^2
\end{equation}
Therefore, we can bound $\delta$ as follows:
    \begin{align}
        \delta = \sum_{|x - E_i| > m} \left\| \sqrt{\frac{\pi}{6m}}\left(\frac{\sin \frac{(x-E_i)\pi}{m}}{ \frac{(x-E_i)\pi}{m}}\right)^m \right\|^2 &\leq \sum_{|x - E_i| > m} \frac{\pi}{6m}\left(\frac{1}{ \frac{(x-E_i)\pi}{m}}\right)^{2m} \\
        &\leq 2\int_{m}^{\infty} \frac{\pi}{6m}\left(\frac{1}{ \frac{x\pi}{m}}\right)^{2m} dx\\
        &\leq  \frac{ \pi^{1-2m}}{3} \frac{1}{2m -1}\\
        &= O\left(\exp (-m)\right)
    \end{align}
\end{proof}
\begin{duplicate}
To achieve a given probability $\delta$ of measuring energy outside of a given confidence interval $\epsilon$ (when using $O(\log \frac{1}{\delta})$-th B-spline boosted QPE), one must use $O(\log \frac{1}{\epsilon} + \log \log \frac{1}{\delta})$ ancillary qubits.
\end{duplicate}
\begin{proof}
Let $l$ be the number of ancilla qubits used by B-spline boosted QPE. Using Lemma \ref{lemma:bspline_prob}, we can achieve a normalised confidence interval of $\frac{2m}{2^l}$ where $m = O(\log \frac{1}{\delta})$. Hence, $\epsilon = \frac{\log \frac{1}{\delta}}{2^{l-1}}$. Therefore,
\begin{equation*}
    l = O\left(\log \left(\frac{1}{\epsilon}\log \frac{1}{\delta}\right)\right)
\end{equation*}
\end{proof}
